\chapter{GPU considerations}
\label{ch:gpu}

The implementation described so far is host code. This chapter states that
plainly, separates what would port readily from what would not, and identifies
the one decision that should be taken before the merge tree is written rather
than after.

\section{Current state}

There are no \texttt{ParallelFor} launches in \texttt{AMRComplex},
\texttt{AMRContour} or \texttt{AMRSegment}. Loops are \texttt{AMREX\_LOOP\_3D},
which expands to serial nested loops on the host, and the working data
structures in the hot paths are \texttt{std::map}, \texttt{std::unordered\_map}
and \texttt{std::set}. None of these is usable in device code as written.

This was a deliberate deferral, not an oversight: correctness on coarse--fine
interfaces and periodic domains was the hard part, and establishing it in host
code kept the verification tractable. It is nonetheless a real limitation for a
code base whose production path is GPU.

\section{What ports readily}

\begin{itemize}[leftmargin=*]
\item \textbf{Node classification} (\cref{sec:identity,sec:stitch}). Index
  arithmetic over boxes with no data-dependent control flow. The
  \texttt{AMREX\_LOOP\_3D} bodies map to \texttt{ParallelFor} essentially
  unchanged; \texttt{Node} is already trivially copyable by construction, so it
  can live in device memory.
\item \textbf{Global numbering} (\cref{eq:gid}). A pure function of
  $(\lev,\iv)$ with no lookup and no communication---evaluable per cell on the
  device.
\item \textbf{Element generation} (\cref{sec:ownership}). Ownership is box
  algebra resolved once per grid on the host; the per-cell corner gather and
  validity test that follow are a device loop. A count-then-fill pass with a
  prefix sum replaces \texttt{push\_back}.
\item \textbf{Simplex subdivision and marching} (\cref{sec:cone}). Per-element,
  branch-light, fixed-size output. Well suited to the device.
\item \textbf{Volume and area reductions.} Already reductions; the compensated
  summation of \cref{ch:verification} needs replacing by a device-friendly
  equivalent, and the integer-derived volumes of \cref{eq:volume} are
  unaffected.
\end{itemize}

\section{What requires redesign}

\begin{itemize}[leftmargin=*]
\item \textbf{The vertex cache} (\cref{sec:edgekey}). A \texttt{std::map} from
  edge key to vertex index, with insertion during traversal. The device
  equivalent is to emit one candidate vertex per crossed edge, sort by the
  $64$-bit edge key, and compact---replacing the associative container by a sort
  and a unique. The edge key is already a pair of integers, so this is a change
  of mechanism rather than of design, and \cref{rule:edgefn} is what makes it
  legitimate: because a vertex is a function of its edge alone, duplicates are
  bitwise identical and compaction is exact.
\item \textbf{Union--find} (\cref{ch:segment}). Pointer-chasing with path
  compression over a hash map. A device implementation needs a different
  algorithm---label propagation over a fixed adjacency, or a device-side
  connected-components pass---together with the adjacency materialised as arrays
  rather than discovered by traversal.
\item \textbf{The face and edge censuses} (\cref{sec:faces,sec:orient}). Both
  build \texttt{std::map} tallies. The exclusive-or parities are already
  reductions and port directly; the exact incidence counts would become
  sort-and-count.
\item \textbf{MEF output} (\cref{ch:contour}). Host-side by nature. The
  \texttt{.pvtu} path on the roadmap is the better target and is equally
  host-side, which is appropriate for I/O.
\end{itemize}

\section{A decision to take before the merge tree}

\decide{Whether the merge tree is written as host code, accepting a
device-to-host transfer of the node field per snapshot, or designed for the
device from the start.}

The consideration is that the merge tree is the first component whose cost scales
with the \emph{volume} of data rather than with the area of an extracted surface.
Contouring touches only cells straddling one isovalue; a sweep does so once per
threshold; the merge tree visits every node in sorted order. A global sort of the
node field is the dominant term, and that is a well-suited device operation
whereas the union--find that consumes it is not.

The pragmatic shape, if the device route is chosen, is a device sort feeding a
host tree construction on the resulting order, since the tree itself is small
compared with the field. That should be settled deliberately rather than
inherited from the host-only precedent set here.

A second point worth recording: the verification apparatus of
\cref{ch:verification} is dimension- and backend-agnostic by construction. The
checksums, the analytic references and the negative controls all compare
\emph{outputs}, so a device port can be gated on reproducing the host results
exactly---which is the strongest possible acceptance test for such a port, and is
available at no additional cost.
